% !TEX root = ../main.tex
%% Corrigé du devoir surveillé n° 2 — valeurs vérifiées : verifications/devoirs.py

\corrigeDevoir{2}

\begin{corrigeBac}{8}{L'eau oxygénée à 10 volumes}
\begin{questions}
\item $\ce{MnO4- + 8 H+ + 5 e- <=> Mn^{2+} + 4 H2O}$ et $\ce{H2O2 <=> O2 + 2 H+ + 2 e-}$. En multipliant la première par 2 et la seconde
par 5, les 10 électrons s'éliminent : $\ce{2 MnO4- + 16 H+ + 5 H2O2 -> 2 Mn^{2+} + 8 H2O + 5 O2 + 10 H+}$, soit, après simplification
des ions $\ce{H+}$, l'\res{\text{équation annoncée}}. \pts{1,5 pt}
\item Burette graduée contenant la solution de permanganate ; erlenmeyer contenant $V_1=\qty{10,0}{mL}$ de $\mathrm{S_1}$ prélevés à la pipette
jaugée ; agitateur magnétique et barreau aimanté (schéma ci-dessous). \pts{1 pt}
\item Avant l'équivalence, les ions $\ce{MnO4-}$ versés sont consommés : la solution reste incolore. L'équivalence est repérée par la
\res{\text{persistance de la teinte violette}} : l'ion permanganate, coloré, joue lui-même le rôle d'indicateur. \pts{1 pt}
\item À l'équivalence, les réactifs ont été introduits dans les \res{\text{proportions stœchiométriques}} :
$\dfrac{n(\ce{H2O2})}{5}=\dfrac{n_\mathrm{E}(\ce{MnO4-})}{2}$, soit $\dfrac{c_1\,V_1}{5}=\dfrac{c_2\,V_\mathrm{E}}{2}$. \pts{1 pt}
\item $c_1=\dfrac{5\,c_2\,V_\mathrm{E}}{2\,V_1}=\dfrac{5\times\num{2,00e-2}\times17{,}8}{2\times10{,}0}=\res{\qty{0,0890}{mol.L^{-1}}}$ ; solution commerciale :
$c_0=10\,c_1=\res{\qty{0,890}{mol.L^{-1}}}$. \pts{1,5 pt}
\item Un litre libère $\dfrac{10}{22{,}4}=\qty{0,446}{mol}$ de dioxygène, ce qui consomme $2\times0{,}446=\qty{0,893}{mol}$ de $\ce{H2O2}$ : la solution
doit avoir une concentration de \qty{0,893}{mol.L^{-1}}. L'écart avec \qty{0,890}{mol.L^{-1}} est de \qty{0,3}{\percent} seulement :
l'indication est \res{\text{vérifiée}}. \pts{1,5 pt}
\item Sans dilution, il faudrait verser $10\times17{,}8=\qty{178}{mL}$ : la \res{\text{burette de 25 mL}} ne suffirait pas. \pts{0,5 pt}
\end{questions}
\schema[corrige]{ds2-montage}
\end{corrigeBac}

\begin{corrigeBac}{6}{Trois molécules}
\begin{questions}
\item Configurations : C $1s^2\,2s^2\,2p^2$ ; N $1s^2\,2s^2\,2p^3$ ; O $1s^2\,2s^2\,2p^4$ ; Cl $1s^2\,2s^2\,2p^6\,3s^2\,3p^5$. Électrons de la
couche externe, ou de valence : C \res{4}, N \res{5}, O \res{6}, Cl \res{7}. \pts{1 pt}
\item $\ce{CH2O}$ : $4+2+6=12$ électrons, 6 doublets ; $\ce{HCN}$ : $1+4+5=10$ électrons, 5 doublets ; $\ce{CCl4}$ : $4+4\times7=32$ électrons,
16 doublets. Schémas ci-dessous : le carbone forme quatre liaisons (une double dans $\ce{CH2O}$, une triple dans $\ce{HCN}$). \pts{1,5 pt}
\item Autour du carbone : trois groupes de doublets dans $\ce{CH2O}$, molécule \res{\text{plane}} (angles de \ang{120}) ; deux dans $\ce{HCN}$,
molécule \res{\text{linéaire}} ; quatre dans $\ce{CCl4}$, molécule \res{\text{tétraédrique}}. \pts{1 pt}
\item Écarts : C$=$O : 0,89 ; C$\equiv$N : 0,49 ; C--Cl : 0,61, liaisons \res{\text{polarisées}}, avec $\delta^-$ sur O, N et Cl et $\delta^+$ sur C.
C--H : 0,35, liaison non polarisée. \pts{1 pt}
\item $\ce{CH2O}$ et $\ce{HCN}$ sont \res{\text{polaires}} : les centres des charges partielles ne coïncident pas. Dans $\ce{CCl4}$, les quatre liaisons
polarisées sont disposées symétriquement : le centre des charges $\delta^-$ coïncide avec l'atome de carbone, la molécule est
\res{\text{apolaire}}. \pts{1 pt}
\item L'eau est polaire : elle dissout bien les espèces polaires, comme le méthanal, qui établit en plus des \res{\text{liaisons
hydrogène}} avec l'eau ; le tétrachlorométhane, apolaire, y est très peu soluble. \pts{0,5 pt}
\end{questions}
\schema[corrige]{ds2-lewis}
\piege{Conclure qu'une molécule est polaire dès qu'elle contient une liaison polarisée : il faut aussi examiner sa géométrie.}
\end{corrigeBac}

\begin{corrigeBac}{6}{Le limonène de l'orange}
\begin{questions}
\item L'eau est un solvant \res{\text{polaire}}, dont les molécules sont liées par des liaisons hydrogène ; le limonène, \res{\text{apolaire}}, ne
peut pas établir d'interactions fortes avec elles. \pts{1 pt}
\item L'\res{\text{éthanol}} : il est miscible à l'eau, on n'obtiendrait pas deux phases à séparer. \pts{1 pt}
\item Les deux solvants conviennent chimiquement ; on préfère le \res{\text{cyclohexane}}, moins dangereux pour la santé que le
dichlorométhane, suspecté d'être cancérogène. \pts{1 pt}
\item On verse le distillat et le cyclohexane dans l'ampoule, on bouche, on agite en dégazant régulièrement, puis on laisse décanter.
Le cyclohexane ($d=0{,}78$) forme la \res{\text{phase supérieure}}, qui contient le limonène ; l'eau forme la \res{\text{phase inférieure}}. \pts{2 pts}
\item On ôte le bouchon, on élimine la phase aqueuse par le robinet, puis on recueille la phase organique dans un bécher. On
\res{\text{évapore le cyclohexane}}, beaucoup plus volatil (\qty{81}{\degreeCelsius}) que le limonène (\qty{176}{\degreeCelsius}). \pts{1 pt}
\end{questions}
\schema[corrige]{ds2-ampoule}
\end{corrigeBac}
